%  LaTeX support: latex@mdpi.com 
%  For support, please attach all files needed for compiling as well as the log file, and specify your operating system, LaTeX version, and LaTeX editor.

%=================================================================
\documentclass[journal,article,submit,pdftex,moreauthors]{Definitions/mdpi} 
%\documentclass[preprints,article,submit,pdftex,moreauthors]{Definitions/mdpi} 
% For posting an early version of this manuscript as a preprint, you may use "preprints" as the journal. Changing "submit" to "accept" before posting will remove line numbers.

%--------------------
% Class Options:
%--------------------
%----------
% journal
%----------
% Choose between the following MDPI journals:
% accountaudit, acoustics, actuators, addictions, adhesives, adolescents, aerobiology, aerospace, agriculture, agriengineering, agrochemicals, agronomy, ai, aichem, aieng, aimater, aimed, aipa, air, aisens, algorithms, allergies, alloys, amh, analog, analytica, analytics, anatomia, anesthres, animals, antibiotics, antibodies, antioxidants, applbiosci, appliedchem, appliedmath, appliedphys, applmech, applmicrobiol, applnano, applsci, aquacj, architecture, arm, arthropoda, asc, asi, astronautics, astronomy, atmosphere, atoms, audiolres, automation, axioms, bacteria, batteries, bdcc, beverages, biochem, bioengineering, biologics, biology, biomass, biomechanics, biomed, biomedicines, biomedinformatics, biomimetics, biomolecules, biophysica, bioresourbioprod, biosensors, biosphere, biotech, birds, blockchains, bloods, blsf, brainsci, breath, buildings, cancers, carbon, cardio, cardiogenetics, cardiovascmed, catalysts, cells, ceramics, challenges, chemengineering, chemistry, chemosensors, chemproc, children, chips, cimb, civileng, cleantechnol, climate, clinbioenerg, clinpract, clockssleep, cmd, cmtr, coasts, coatings, colloids, colorants, commodities, complexities, complications, compounds, computation, computers, condensedmatter, conservation, constrmater, cosmetics, covid, crops, cryo, cryptography, crystals, csmf, ctn, culture, curroncol, cyber, dairy, data, ddc, dentistry, dermato, dermatopathology, designs, devices, dhi, diabetology, diagnostics, dietetics, digital, disabilities, diseases, diversity, dna, drones, dynamics, earth, ebj, ecm, ecologies, edm, eesp, electricity, electrochem, electronicmat, electronics, encyclopedia, endocrines, energies, eng, engproc, entomology, entropy, environments, environremediat, epidemiologia, epigenomes, esa, est, fermentation, fibers, fintech, fire, fishes, fluids, foods, forecasting, forensicsci, forests, fossstud, foundations, fractalfract, fuels, future, futureinternet, futurepharmacol, futurephys, futuretransp, galaxies, gases, gastroent, gastrointestdisord, gastronomy, gels, genes, geographies, geohazards, geomatics, geometry, geosciences, geotechnics, geriatrics, germs, glacies, grasses, green, greenhealth, gucdd, hardware, hazardousmatters, healthcare, hearts, hemato, hematolrep, hep, heritage, higheredu, highthroughput, horticulturae, hospitals, hydrobiology, hydrogen, hydrology, hydropower, hygiene, idr, iic, ijem, ijerph, ijgi, ijmd, ijms, ijns, ijom, ijpb, ijt, ijtm, ijtpp, ime, immuno, informatics, information, infrastructures, inorganics, insects, instruments, inventions, iot, j, jaestheticmed, jal, jcdd, jcm, jcp, jcrm, jcs, jcto, jdad, jdb, jdream, jemr, jeta, jfb, jfmk, jgbg, jgg, jimaging, jlpea, jmahp, jmmp, jmms, jmp, jmse, jne, jnt, jof, joi, joitmc, joma, jop, joptm, jor, jox, jpbi, jphytomed, jpm, jsan, jtaer, jvd, jzbg, kidneydial, kinasesphosphatases, knowledge, labmed, laboratories, lae, land, life, lights, limnolrev, lipidology, liquids, livers, logics, logistics, lubricants, lymphatics, machines, macromol, magnetism, magnetochemistry, make, marinedrugs, materials, materproc, mathematics, mca, measurements, medicina, medicines, medsci, membranes, merits, metabolites, metals, meteorology, methane, metrics, metrology, micro, microarrays, microbiolres, microelectronics, micromachines, microorganisms, microplastics, microwave, minerals, mining, mmphys, modelling, molbank, molecules, mps, msf, mti, multimedia, muscles, nanoenergyadv, nanomanufacturing, nanomaterials, ncrna, ndt, network, neuroglia, neuroimaging, neurolint, neurosci, nitrogen, notspecified, nursrep, nutraceuticals, nutrients, obesities, occuphealth, oceans, ohbm, onco, optics, oral, organics, organoids, osteology, oxygen, pandemics, parasites, parasitologia, particles, pathogens, pathophysiology, pediatrrep, pets, pharmaceuticals, pharmaceutics, pharmacoepidemiology, pharmacy, philosophies, photochem, photonics, phycology, physchem, physics, physiologia, plants, plasma, platforms, pollutants, polymers, polysaccharides, populations, poultry, powders, precipitation, precisoncol, preprints, proceedings, processes, prosthesis, proteomes, psf, psychiatryint, psychoactives, purification, quantumrep, quaternary, qubs, radiation, rdt, reactions, realestate, receptors, recycling, regeneration, remotesensing, reports, reprodmed, resources, rheumato, rjpm, robotics, rsee, ruminants, safety, sci, scipharm, sclerosis, seeds, sensors, separations, sexes, shi, signals, sinusitis, siuj, skins, smartcities, sna, societies, software, soilsystems, solar, solids, spectroscj, sports, standards, stats, std, stratsediment, stresses, surfaces, surgeries, suschem, sustainability, symmetry, synbio, systems, tae, targets, taxonomy, technologies, telecom, test, textiles, thalassrep, therapeutics, thermo, timespace, tomography, toxics, toxins, tph, transplantology, transportation, traumacare, traumas, tri, tropicalmed, universe, urbansci, uro, vaccines, vehicles, venereology, vetsci, vibration, virtualworlds, viruses, vision, waste, water, welding, wem, wevj, wild, wind, women, world, zoonoticdis

%---------
% article
%---------
% The default type of manuscript is "article", but can be replaced by: 
% abstract, addendum, article, benchmark, book, bookreview, briefcommunication, briefreport, casereport, changes, clinicopathologicalchallenge, comment, commentary, communication, conceptpaper, conferenceproceedings, correction, conferencereport, creative, datadescriptor, discussion, entry, expressionofconcern, extendedabstract, editorial, essay, erratum, fieldguide, hypothesis, interestingimages, letter, meetingreport, monograph, newbookreceived, obituary, opinion, proceedingpaper, projectreport, reply, retraction, review, perspective, protocol, shortnote, studyprotocol, supfile, systematicreview, technicalnote, viewpoint, guidelines, registeredreport, tutorial,  giantsinurology, urologyaroundtheworld
% supfile = supplementary materials

%----------
% submit
%----------
% The class option "submit" will be changed to "accept" by the Editorial Office when the paper is accepted. This will only make changes to the frontpage (e.g., the logo of the journal will get visible), the headings, and the copyright information. Also, line numbering will be removed. Journal info and pagination for accepted papers will also be assigned by the Editorial Office.

%------------------
% moreauthors
%------------------
% If there is only one author the class option oneauthor should be used. Otherwise use the class option moreauthors.

%---------
% pdftex
%---------
% The option pdftex is for use with pdfLaTeX. Remove "pdftex" for (1) compiling with LaTeX & dvi2pdf (if eps figures are used) or for (2) compiling with XeLaTeX.

%=================================================================
% MDPI internal commands - do not modify
\firstpage{1} 
\makeatletter 
\setcounter{page}{\@firstpage} 
\makeatother
\pubvolume{1}
\issuenum{1}
\articlenumber{0}
\pubyear{2026}
\copyrightyear{2026}
%\externaleditor{Firstname Lastname} % More than 1 editor, please add `` and '' before the last editor name
\datereceived{ } 
\daterevised{ } % Comment out if no revised date
\dateaccepted{ } 
\datepublished{ } 
%\datecorrected{} % For corrected papers: "Corrected: XXX" date in the original paper.
%\dateretracted{} % For retracted papers: "Retracted: XXX" date in the original paper.
%\doinum{} % Used for some special journals, like molbank
%\pdfoutput=1 % Uncommented for upload to arXiv.org
%\CorrStatement{yes}  % For updates
%\longauthorlist{yes} % For many authors that exceed the left citation part
%\IsAssociation{yes} % For association journals

%=================================================================
% Add packages and commands here. The following packages are loaded in our class file: fontenc, inputenc, calc, indentfirst, fancyhdr, graphicx, epstopdf, lastpage, ifthen, float, amsmath, amssymb, lineno, setspace, enumitem, mathpazo, booktabs, titlesec, etoolbox, tabto, xcolor, colortbl, soul, multirow, microtype, tikz, totcount, changepage, attrib, upgreek, array, tabularx, pbox, ragged2e, tocloft, marginnote, marginfix, enotez, amsthm, natbib, hyperref, cleveref, scrextend, url, geometry, newfloat, caption, draftwatermark, seqsplit
% cleveref: load \crefname definitions after \begin{document}

\usepackage{pgfplots}
\pgfplotsset{compat=1.18}
\definecolor{regularColor}{HTML}{4C78A8}
\definecolor{codeColor}{HTML}{F58518}

%=================================================================
% Please use the following mathematics environments: Theorem, Lemma, Corollary, Proposition, Characterization, Property, Problem, Example, ExamplesandDefinitions, Hypothesis, Remark, Definition, Notation, Assumption
%% For proofs, please use the proof environment (the amsthm package is loaded by the MDPI class).

%=================================================================
% Full title of the paper (Capitalized)
\Title{Code-Generated Tool Orchestration versus Native Function Calling in Stateful Business Workflows: A Multi-Model Benchmark Study}

% Author Orchid ID: enter ID or remove command
\newcommand{\orcidauthorA}{0000-0000-0000-000X} % Update before submission
%\newcommand{\orcidauthorB}{0000-0000-0000-000X} % Add \orcidB{} behind the author's name

% Authors, for the paper (add full first names)
\Author{Anonymous Author $^{1}$\orcidA{} and Anonymous Author $^{1,}$*}

%\longauthorlist{yes}

% MDPI internal command: Authors, for metadata in PDF
\AuthorNames{Anonymous Author and Anonymous Author}

% Affiliations / Addresses (Add [1] after \address if there is only one affiliation.)
\address{%
$^{1}$ \quad Department of Computer Science, Institution Name, City, Country; author@institution.edu}

% Contact information of the corresponding author
\corres{Correspondence: author@institution.edu}

% Current address and/or shared authorship
%\firstnote{Current address: Affiliation.}  
% Current address should not be the same as any items in the Affiliation section.

%\secondnote{These authors contributed equally to this work.}
% The commands \thirdnote{} till \eighthnote{} are available for further notes.

%\simplesumm{} % Simple summary

%\conference{} % An extended version of a conference paper

% Abstract (Do not insert blank lines, i.e. \\) 
\abstract{Large language model agents are commonly deployed with native function calling, but a competing design pattern asks the model to generate executable code that performs multiple tool calls inside a sandbox. We present a reproducible benchmark that compares these two agent architectures on eight stateful business workflow scenarios with strict end-state validation. The benchmark includes multi-provider model routing, progressive tool discovery over a virtual filesystem interface, deterministic scenario validation, and execution-level observability traces. We evaluate seven model configurations across OpenAI-, Anthropic-, and OpenAI-compatible endpoints. Results show that performance depends strongly on model family: for GPT-5.3 Codex and Claude Opus 4.6, code-generated orchestration reduces total token use substantially while preserving perfect validation; for several OpenRouter-routed models, code mode exhibits high variance, increased retries, and occasional validation regressions. Across all models, code mode offers an architectural advantage when long tool chains can be batched into one successful code execution, but this advantage degrades under repeated repair loops and unstable backend behavior. Our findings indicate that ``code mode'' should be treated as a model- and provider-conditional optimization rather than a universal default. We release a practical evaluation framework and a detailed error taxonomy to support production routing decisions for tool-using agents.}

% Keywords
\keyword{large language model agents; tool calling; code generation; sandbox execution; model context protocol; benchmarking; observability; stateful workflows} 

% The fields PACS, MSC, and JEL may be left empty or commented out if not applicable
%\PACS{J0101}
%\MSC{}
%\JEL{}

%%%%%%%%%%%%%%%%%%%%%%%%%%%%%%%%%%%%%%%%%%
% Only for the journal Diversity
%\LSID{\url{http://}}

%%%%%%%%%%%%%%%%%%%%%%%%%%%%%%%%%%%%%%%%%%
% Only for the journal Applied Sciences
%\featuredapplication{Authors are encouraged to provide a concise description of the specific application or a potential application of the work. This section is not mandatory.}
%%%%%%%%%%%%%%%%%%%%%%%%%%%%%%%%%%%%%%%%%%

%%%%%%%%%%%%%%%%%%%%%%%%%%%%%%%%%%%%%%%%%%
% Only for the journal Data
%\dataset{DOI number or link to the deposited data set if the data set is published separately. If the data set shall be published as a supplement to this paper, this field will be filled by the journal editors. In this case, please submit the data set as a supplement.}
%\datasetlicense{License under which the data set is made available (CC0, CC-BY, CC-BY-SA, CC-BY-NC, etc.)}

%%%%%%%%%%%%%%%%%%%%%%%%%%%%%%%%%%%%%%%%%%
% Only for the journal BioTech, Fishes, Neuroimaging and Toxins
%\keycontribution{The breakthroughs or highlights of the manuscript. Authors can write one or two sentences to describe the most important part of the paper.}

%%%%%%%%%%%%%%%%%%%%%%%%%%%%%%%%%%%%%%%%%%
% Only for the journal Encyclopedia
%\encyclopediadef{For entry manuscripts only: please provide a brief overview of the entry title instead of an abstract.}

%%%%%%%%%%%%%%%%%%%%%%%%%%%%%%%%%%%%%%%%%%
% Different journals have different requirements. Please check the specific journal guidelines in the "Instructions for Authors" on the journal's official website.
%\addhighlights{yes}
%\renewcommand{\addhighlights}{%
%
%\noindent The goal is to increase the discoverability and readability of the article via search engines and other scholars. Highlights should not be a copy of the abstract, but a simple text allowing the reader to quickly and simplified find out what the article is about and what can be cited from it. Each of these parts should be devoted up to 2~bullet points.\vspace{3pt}\\
%\textbf{What are the main findings?}
% \begin{itemize}[labelsep=2.5mm,topsep=-3pt]
% \item First bullet.
% \item Second bullet.
% \end{itemize}\vspace{3pt}
%\textbf{What are the implications of the main findings?}
% \begin{itemize}[labelsep=2.5mm,topsep=-3pt]
% \item First bullet.
% \item Second bullet.
% \end{itemize}
%}

%%%%%%%%%%%%%%%%%%%%%%%%%%%%%%%%%%%%%%%%%%
\begin{document}

%%%%%%%%%%%%%%%%%%%%%%%%%%%%%%%%%%%%%%%%%%
\section{Introduction}

Tool-using large language model (LLM) agents are increasingly deployed for enterprise automation. Most production systems still use native function calling loops, where the model emits one tool call per reasoning turn and receives intermediate results back into the next prompt turn. An alternative architecture, here called \emph{code-generated orchestration} (or \emph{code mode}), asks the model to generate executable code that can invoke many tools within one sandboxed run. This architectural change can alter latency, token consumption, and failure dynamics.

This paper presents a full-stack benchmark for comparing these two architectures under controlled, stateful workloads. Unlike stateless toy tasks, our scenarios mutate shared accounting and operations state, requiring ordered multi-step execution and measurable final-state correctness. We focus on the following research questions:

\begin{itemize}
\item \textbf{RQ1}: Under equivalent tools and scenarios, does code-generated orchestration improve end-to-end efficiency (time, tokens, iterations)?
\item \textbf{RQ2}: How does the effect vary across model families and provider routing paths?
\item \textbf{RQ3}: What failure patterns dominate each architecture, and how observable are they?
\item \textbf{RQ4}: When should production systems prefer regular tool calling versus code mode?
\end{itemize}

The benchmark includes seven model configurations, eight realistic business scenarios, strict validation rules, and granular observability traces. A key contribution is the operational perspective: we do not claim a universal winner. Instead, we characterize the boundary conditions under which code mode is a strong optimization versus a regression risk.

%%%%%%%%%%%%%%%%%%%%%%%%%%%%%%%%%%%%%%%%%%
\section{Materials and Methods}

\subsection{System Overview}

The benchmark compares two agent classes against the same tool backend and scenario set:

\begin{itemize}
\item \textbf{Regular agent}: iterative, native tool/function calling via provider APIs.
\item \textbf{Code mode agent}: model-generated Python code executed in a restricted sandbox with tool bindings.
\end{itemize}

Each scenario run resets global state before execution, ensuring independence. Both agent outputs are validated against scenario-specific predicates defined over final state.

\subsection{Deep Implementation Internals}

\textbf{Agent routing and runtime resolution.} The agent factory maps model keys to provider adapters, model names, and fallback routing logic. Several models support direct-provider override (for example, direct Zhipu or MiniMax endpoints) when dedicated keys are available; otherwise, requests route through OpenRouter-compatible endpoints.

\textbf{Regular agent loop.} The regular path follows a classical loop: prompt with schema, receive tool call(s), execute tool(s), return tool results, and continue until final textual response or max-iteration termination.

\textbf{Code mode loop.} The code-mode path prompts for a single Python block, executes it in a restricted environment, and checks whether variable \texttt{result} is produced. On failure, a retry prompt injects compact error hints and constraints. State snapshot/restore is used to roll back failed attempts.

\textbf{Sandboxing and execution controls.} Generated code is compiled with RestrictedPython, guarded with:

\begin{itemize}
\item import allow-list,
\item timeout budget,
\item memory cap,
\item restricted builtins,
\item tool-call interception and audit log capture.
\end{itemize}

\textbf{Progressive tool discovery.} In addition to direct tool method calls, code mode supports virtual filesystem navigation with \texttt{tools.ls(path)}, \texttt{tools.read(path)}, and \texttt{tools.call(path,args)}. This enables API discovery without forcing all schemas into every prompt context.

\textbf{Observability.} For code mode, each execution attempt records:

\begin{itemize}
\item model response text,
\item extracted code,
\item sandbox metrics,
\item per-call tool arguments and outputs,
\item state deltas,
\item expected-vs-actual tool-flow discrepancies.
\end{itemize}

\subsection{Workload Design}

We use eight scenarios spanning expense recording, multi-client invoicing, reconciliation, account transfer management, quarter-end synthesis, and category-level budget analysis. Scenarios encode expected minimum tool usage and expected terminal-state checks (for example invoice statuses, balances, transaction counts, and receivables).

\subsection{Evaluation Metrics}

For each model and architecture pair, we report:

\begin{itemize}
\item validation pass rate,
\item average scenario runtime,
\item average iteration count,
\item total input and output tokens,
\item code-mode iteration failures,
\item code-mode tool discrepancies.
\end{itemize}

Token-based cost can be estimated as:
\begin{linenomath}
\begin{equation}
\mathrm{cost}=\frac{T_{in}}{10^6}R_{in}+\frac{T_{out}}{10^6}R_{out},
\end{equation}
\end{linenomath}
where $T_{in}$ and $T_{out}$ are input/output token totals and $R_{in}$, $R_{out}$ are provider rates.

%%%%%%%%%%%%%%%%%%%%%%%%%%%%%%%%%%%%%%%%%%
\section{Results}

\subsection{Model-Level Aggregate Outcomes}

Table~\ref{tab:model_results} summarizes model-level outcomes from March 2026 benchmark artifacts. To avoid bias from incomplete execution, the primary quantitative comparison uses six complete model runs. Kimi 2.5 is retained as a pending-completion entry and excluded from comparative figures.

\begin{table}[H]
\caption{Aggregate benchmark outcomes by model. Time is mean scenario runtime. Tokens are total in+out over successful scenarios.\label{tab:model_results}}
\begin{tabularx}{\textwidth}{>{\RaggedRight\arraybackslash}XCCCC}
\toprule
\textbf{Model} & \textbf{Validation (R/C)} & \textbf{Time (R/C)} & \textbf{Iterations (R/C)} & \textbf{Tokens (R/C)}\\
\midrule
GPT-5.3 Codex & 8/8 vs 8/8 & 17.77s vs 14.86s & 5.25 vs 1.25 & 95,482 vs 22,352\\
GPT-5.2 & 8/8 vs 6/8 & 12.33s vs 27.23s & 4.12 vs 2.75 & 79,240 vs 79,657\\
Claude Opus 4.6 & 8/8 vs 8/8 & 27.30s vs 19.53s & 4.12 vs 1.38 & 435,483 vs 27,292\\
GLM-5 (OpenRouter route) & 7/8 vs 7/8 & 77.19s vs 213.53s & 4.00 vs 2.62 & 105,037 vs 92,367\\
MiniMax M2.5 & 7/8 vs 5/8 & 122.43s vs 53.07s & 4.12 vs 7.00 & 125,733 vs 104,001\\
Gemini 3.1 Flash Lite (route) & 8/8 vs 4/8 & 8.54s vs 4.95s & 4.00 vs 6.50 & 79,394 vs 26,003\\
Kimi 2.5 (partial run) & 6/7 vs 5/6 & mixed & mixed & not stable in summary\\
\bottomrule
\end{tabularx}
\end{table}

\subsection{Validation, Latency, and Token Efficiency}

The validation comparison (Figure~\ref{fig:validation_rate}) shows a bifurcated pattern. Two high-capability native-path models (GPT-5.3 Codex and Claude Opus 4.6) preserve full validation under both architectures. In contrast, some routed or smaller configurations experience validation degradation under code mode (for example Gemini 3.1 Flash Lite route: $100\%\rightarrow 50\%$; MiniMax M2.5: $87.5\%\rightarrow 62.5\%$).

\begin{figure}[H]
\centering
\begin{tikzpicture}
\begin{axis}[
	width=\textwidth,
	height=6.2cm,
	ybar,
	bar width=8pt,
	ymin=0,
	ymax=110,
	ylabel={Validation Pass Rate (\%)},
	symbolic x coords={GPT-5.3,GPT-5.2,Opus4.6,GLM-5,MiniMax,Gemini3.1},
	xtick=data,
	x tick label style={rotate=20,anchor=east,font=\footnotesize},
	legend style={at={(0.5,1.05)},anchor=south,legend columns=2},
	ymajorgrids=true,
]
\addplot[fill=regularColor] coordinates {(GPT-5.3,100) (GPT-5.2,100) (Opus4.6,100) (GLM-5,87.5) (MiniMax,87.5) (Gemini3.1,100)};
\addplot[fill=codeColor] coordinates {(GPT-5.3,100) (GPT-5.2,75) (Opus4.6,100) (GLM-5,87.5) (MiniMax,62.5) (Gemini3.1,50)};
\legend{Regular,Code Mode}
\end{axis}
\end{tikzpicture}
\caption{Validation pass rate comparison by model (complete runs only).\label{fig:validation_rate}}
\end{figure}

Latency behavior (Figure~\ref{fig:latency_compare}) is architecture- and model-dependent rather than monotonic. Code mode is faster for GPT-5.3 Codex ($-16.4\%$), Opus 4.6 ($-28.5\%$), MiniMax M2.5 ($-56.6\%$), and Gemini 3.1 route ($-42.0\%$), but slower for GPT-5.2 ($+120.8\%$) and GLM-5 route ($+176.6\%$).

\begin{figure}[H]
\centering
\begin{tikzpicture}
\begin{axis}[
	width=\textwidth,
	height=6.2cm,
	ybar,
	bar width=8pt,
	ymin=0,
	ymax=240,
	ylabel={Average Runtime (s)},
	symbolic x coords={GPT-5.3,GPT-5.2,Opus4.6,GLM-5,MiniMax,Gemini3.1},
	xtick=data,
	x tick label style={rotate=20,anchor=east,font=\footnotesize},
	legend style={at={(0.5,1.05)},anchor=south,legend columns=2},
	ymajorgrids=true,
]
\addplot[fill=regularColor] coordinates {(GPT-5.3,17.77) (GPT-5.2,12.33) (Opus4.6,27.30) (GLM-5,77.19) (MiniMax,122.43) (Gemini3.1,8.54)};
\addplot[fill=codeColor] coordinates {(GPT-5.3,14.86) (GPT-5.2,27.23) (Opus4.6,19.53) (GLM-5,213.53) (MiniMax,53.07) (Gemini3.1,4.95)};
\legend{Regular,Code Mode}
\end{axis}
\end{tikzpicture}
\caption{Average end-to-end runtime by model and architecture.\label{fig:latency_compare}}
\end{figure}

Token behavior (Figure~\ref{fig:token_compare}) reveals the strongest and most consistent upside of code mode for robust code-generating models. GPT-5.3 Codex reduces aggregate tokens by approximately $76.6\%$, while Opus 4.6 reduces by approximately $93.7\%$. Conversely, GPT-5.2 exhibits near parity in aggregate tokens and significantly higher latency due to retries.

\begin{figure}[H]
\centering
\begin{tikzpicture}
\begin{axis}[
	width=\textwidth,
	height=6.2cm,
	ybar,
	bar width=8pt,
	ymode=log,
	log basis y=10,
	ymin=10000,
	ymax=1000000,
	ylabel={Total Tokens (log scale)},
	symbolic x coords={GPT-5.3,GPT-5.2,Opus4.6,GLM-5,MiniMax,Gemini3.1},
	xtick=data,
	x tick label style={rotate=20,anchor=east,font=\footnotesize},
	legend style={at={(0.5,1.05)},anchor=south,legend columns=2},
	ymajorgrids=true,
]
\addplot[fill=regularColor] coordinates {(GPT-5.3,95482) (GPT-5.2,79240) (Opus4.6,435483) (GLM-5,105037) (MiniMax,125733) (Gemini3.1,79394)};
\addplot[fill=codeColor] coordinates {(GPT-5.3,22352) (GPT-5.2,79657) (Opus4.6,27292) (GLM-5,92367) (MiniMax,104001) (Gemini3.1,26003)};
\legend{Regular,Code Mode}
\end{axis}
\end{tikzpicture}
\caption{Total token consumption by model and architecture (log scale).\label{fig:token_compare}}
\end{figure}

\subsection{Primary Patterns}

Three strong patterns emerge:

\begin{enumerate}
\item \textbf{Code mode is highly effective for some high-capability models.} GPT-5.3 Codex and Claude Opus 4.6 retain full validation while sharply reducing total tokens.
\item \textbf{Code mode can regress under repeated repair loops.} GPT-5.2 and several routed models show slower total runtime due to multi-iteration code repair.
\item \textbf{Provider path influences architecture outcomes.} OpenAI/Anthropic-native results are more stable than several OpenRouter-routed runs in this dataset.
\end{enumerate}

\subsection{Failure and Observability Analysis}

Code-mode observability traces identify two dominant failure classes:

\begin{itemize}
\item \textbf{Execution-contract failures}: generated code fails sandbox constraints or misses required \texttt{result} assignment.
\item \textbf{Tool-sequence mismatches}: generated workflows complete partially but violate expected operation ordering or counts.
\end{itemize}

A practical finding is that low-level sandbox overhead is usually tiny (millisecond scale) relative to LLM round-trip time; total runtime is governed primarily by number of model iterations, not execution engine latency. For example, in several failing high-latency runs, sandbox compile and execution remain sub-100 ms while overall wall-clock runtime expands due to repeated model retries.

\subsection{Answer to Research Questions}

\textbf{RQ1 (efficiency).} Code mode can substantially reduce token footprint and latency, but only when first-pass or low-retry code execution is reliable.

\textbf{RQ2 (model/provider variation).} Variation is large across model-provider combinations; architecture effects are not transferable across all routing paths.

\textbf{RQ3 (failure patterns).} The most costly failures are not sandbox compute failures but iterative regeneration loops triggered by execution contract and tool-flow issues.

\textbf{RQ4 (deployment choice).} A static architecture choice is inferior to policy-based routing that incorporates model reliability history and live retry behavior.

%%%%%%%%%%%%%%%%%%%%%%%%%%%%%%%%%%%%%%%%%%
\section{Discussion}

\subsection{When Code Mode Wins}

Code mode is most compelling when a model can synthesize correct multi-step code in one or two iterations. In that regime, the architecture removes repeated "LLM-think/tool-call/LLM-think" cycles and reduces token amplification from iterative context growth.

\subsection{When Regular Tool Calling Is Safer}

Regular tool calling remains robust under models that struggle with executable code synthesis or repair. The explicit incremental loop can reduce catastrophic divergence in long scenarios, even if token overhead is higher.

\subsection{Operational Routing Implication}

A single global architecture is suboptimal. The data support a policy-driven router:

\begin{itemize}
\item default to code mode for models with demonstrated one-pass reliability,
\item default to regular mode for unstable routed models,
\item dynamically fail over based on early iteration-failure signals.
\end{itemize}

%%%%%%%%%%%%%%%%%%%%%%%%%%%%%%%%%%%%%%%%%%
\section{Limitations and Future Work}

This study has four limitations. First, scenarios are concentrated in business/accounting domains. Second, model versions and provider behavior can drift rapidly over time. Third, cost analysis is token-based and excludes provider-specific queueing effects. Fourth, we do not yet include human preference evaluation for final textual explanations.

Future work should extend domain diversity (support, procurement, planning, and mixed transactional workloads), run repeated seeds with confidence intervals, and evaluate adaptive hybrid agents that begin in code mode and downgrade to regular mode after bounded retries.

%%%%%%%%%%%%%%%%%%%%%%%%%%%%%%%%%%%%%%%%%%
\section{Conclusions}

We introduced a reproducible benchmark that isolates the architectural trade-off between regular native function calling and code-generated tool orchestration. The core conclusion is conditional: code mode is not universally superior, but can be dramatically more efficient for model-provider combinations with strong first-pass code reliability. For unstable combinations, retries erase or reverse gains. Production systems should therefore make architecture a routing decision, informed by validation history, iteration behavior, and provider stability, rather than a static design choice.

%%%%%%%%%%%%%%%%%%%%%%%%%%%%%%%%%%%%%%%%%%
\section{Patents}

No patents are associated with this work.

%% Example of a page in landscape format (with table and table footnote).
%\startlandscape
%\begin{table}[H] %% Table in wide page
%%\isPreprints{\centering}{} % This command is only used for ``preprints''.
%\caption{This is a very wide table.\label{tab3}}
%	\begin{tabularx}{\textwidth}{CCCC}
%		\toprule
%		\textbf{Title 1}	& \textbf{Title 2}	& \textbf{Title 3}	& \textbf{Title 4}\\
%		\midrule
%		Entry 1		& Data			& Data			& This cell has some longer content that runs over two lines.\\
%		Entry 2		& Data			& Data			& Data\textsuperscript{1}\\
%		\bottomrule
%	\end{tabularx}
%%\isPreprints{}{% This command is only used for ``preprints''.
%	\begin{adjustwidth}{+\extralength}{0cm}
%%} % If the paper is ``preprints'', please uncomment this parenthesis.
%		\noindent\footnotesize{\textsuperscript{1} This is a table footnote.}
%%\isPreprints{}{% This command is only used for ``preprints''.
%	\end{adjustwidth}
%%} % If the paper is ``preprints'', please uncomment this parenthesis.
%\end{table}
%\finishlandscape

\vspace{6pt} 

%%%%%%%%%%%%%%%%%%%%%%%%%%%%%%%%%%%%%%%%%%
%% optional
%\supplementary{The following supporting information can be downloaded at:  \linksupplementary{s1}, Figure S1: title; Table S1: title; Video S1: title.}

% Only for journal Methods and Protocols:
% If you wish to submit a video article, please do so with any other supplementary material.
% \supplementary{The following supporting information can be downloaded at: \linksupplementary{s1}, Figure S1: title; Table S1: title; Video S1: title. A supporting video article is available at doi: link.}

% Only used for preprtints:
% \supplementary{The following supporting information can be downloaded at the website of this paper posted on \href{https://www.preprints.org/}{Preprints.org}.}

% Only for journal Hardware:
% If you wish to submit a video article, please do so with any other supplementary material.
% \supplementary{The following supporting information can be downloaded at: \linksupplementary{s1}, Figure S1: title; Table S1: title; Video S1: title.\vspace{6pt}\\
%\begin{tabularx}{\textwidth}{lll}
%\toprule
%\textbf{Name} & \textbf{Type} & \textbf{Description} \\
%\midrule
%S1 & Python script (.py) & Script of python source code used in XX \\
%S2 & Text (.txt) & Script of modelling code used to make Figure X \\
%S3 & Text (.txt) & Raw data from experiment X \\
%S4 & Video (.mp4) & Video demonstrating the hardware in use \\
%... & ... & ... \\
%\bottomrule
%\end{tabularx}
%}

%%%%%%%%%%%%%%%%%%%%%%%%%%%%%%%%%%%%%%%%%%
\authorcontributions{Conceptualization, A.A.; methodology, A.A. and B.B.; software, A.A.; validation, A.A. and B.B.; formal analysis, A.A.; investigation, A.A.; data curation, A.A.; writing---original draft preparation, A.A.; writing---review and editing, A.A. and B.B.; visualization, A.A.; supervision, B.B. All authors have read and agreed to the published version of the manuscript.}

\funding{This research received no external funding.}

\institutionalreview{Not applicable.}

\informedconsent{Not applicable.}

\dataavailability{All benchmark code, scenario definitions, model result JSON files, trace logs, and generated reports are available in the project repository structure used for this study.}

%\durcstatement{Current research is limited to the [please insert a specific academic field, e.g., XXX], which is beneficial [share benefits and/or primary use] and does not pose a threat to public health or national security. Authors acknowledge the dual-use potential of the research involving xxx and confirm that all necessary precautions have been taken to prevent potential misuse. As an ethical responsibility, authors strictly adhere to relevant national and international laws about DURC. Authors advocate for responsible deployment, ethical considerations, regulatory compliance, and transparent reporting to mitigate misuse risks and foster beneficial outcomes.}

% Only for journal Nursing Reports
%\publicinvolvement{Please describe how the public (patients, consumers, carers) were involved in the research. Consider reporting against the GRIPP2 (Guidance for Reporting Involvement of Patients and the Public) checklist. If the public were not involved in any aspect of the research add: ``No public involvement in any aspect of this research''.}
%
%% Only for journal Nursing Reports
%\guidelinesstandards{Please add a statement indicating which reporting guideline was used when drafting the report. For example, ``This manuscript was drafted against the XXX (the full name of reporting guidelines and citation) for XXX (type of research) research''. A complete list of reporting guidelines can be accessed via the equator network: \url{https://www.equator-network.org/}.}
%
%% Only for journal Nursing Reports
%\useofartificialintelligence{Please describe in detail any and all uses of artificial intelligence (AI) or AI-assisted tools used in the preparation of the manuscript. This may include, but is not limited to, language translation, language editing and grammar, or generating text. Alternatively, please state that “AI or AI-assisted tools were not used in drafting any aspect of this manuscript”.}

\acknowledgments{The authors thank internal reviewers for feedback on benchmark design and reporting structure.}

\conflictsofinterest{The authors declare no conflicts of interest.} 

%%%%%%%%%%%%%%%%%%%%%%%%%%%%%%%%%%%%%%%%%%
%% Optional

%% Only for journal Encyclopedia
%\entrylink{The Link to this entry published on the encyclopedia platform.}

\abbreviations{Abbreviations}{
The following abbreviations are used in this manuscript:
\\

\noindent 
\begin{tabular}{@{}ll}
LLM & Large language model\\
MCP & Model Context Protocol\\
VFS & Virtual filesystem\\
RQ & Research question
\end{tabular}
}

%%%%%%%%%%%%%%%%%%%%%%%%%%%%%%%%%%%%%%%%%%
%% Optional
\appendixtitles{no} % Leave argument "no" if all appendix headings stay EMPTY (then no dot is printed after "Appendix A"). If the appendix sections contain a heading then change the argument to "yes".
\appendixstart
\appendix
\section[\appendixname~\thesection]{}
\subsection[\appendixname~\thesubsection]{}
This appendix summarizes benchmark internals needed for exact replication:
\begin{itemize}
\item Scenario source: eight stateful workflows with deterministic validators.
\item Tool backend: shared mutable accounting and operations state.
\item Sandbox limits: timeout and memory boundaries with import restrictions.
\item Artifacts: per-model JSON results, markdown reports, and code-mode traces.
\end{itemize}

\begin{table}[H] 
\caption{Core benchmark artifact locations.\label{tab5}}
%\newcolumntype{C}{>{\centering\arraybackslash}X}
\begin{tabularx}{\textwidth}{CCC}
\toprule
\textbf{Artifact} & \textbf{Path} & \textbf{Purpose}\\
\midrule
Scenarios & tests/test\_scenarios.py & Workload definitions and validators\\
Reports & results/reports/ & Human-readable per-model summaries\\
Traces & results/traces/ & Iteration-level code-mode observability\\
\bottomrule
\end{tabularx}
\end{table}

\section[\appendixname~\thesection]{}
All appendix sections are cited from the main text and provide implementation-level reproducibility details.

%%%%%%%%%%%%%%%%%%%%%%%%%%%%%%%%%%%%%%%%%%
%\isPreprints{}{% This command is only used for ``preprints''.
\begin{adjustwidth}{-\extralength}{0cm}
%} % If the paper is ``preprints'', please uncomment this parenthesis.
%\printendnotes[custom] % Un-comment to print a list of endnotes

\reftitle{References}

% Please provide the correct journal abbreviation (e.g. according to the “List of Title Word Abbreviations” http://www.issn.org/services/online-services/access-to-the-ltwa/).
% Citations and References in Supplementary files are permitted provided that they also appear in the reference list here. 

%=====================================
% References, variant A: external bibliography
%=====================================
% \bibliography{your_external_BibTeX_file}

%=====================================
% References, variant B: internal bibliography
%=====================================

% ACS format
\begin{thebibliography}{999}
\bibitem{refs-pending}
References intentionally omitted at this drafting stage and will be inserted in the citation pass.
\end{thebibliography}

% If authors have biography, please use the format below
%\section*{Short Biography of Authors}
%\bio
%{\raisebox{-0.35cm}{\includegraphics[width=3.5cm,height=5.3cm,clip,keepaspectratio]{Definitions/author1.pdf}}}
%{\textbf{Firstname Lastname} Biography of first author}
%
%\bio
%{\raisebox{-0.35cm}{\includegraphics[width=3.5cm,height=5.3cm,clip,keepaspectratio]{Definitions/author2.jpg}}}
%{\textbf{Firstname Lastname} Biography of second author}

% For the MDPI journals use author-date citation, please follow the formatting guidelines on http://www.mdpi.com/authors/references
% To cite two works by the same author: \citeauthor{ref-journal-1a} (\citeyear{ref-journal-1a}, \citeyear{ref-journal-1b}). This produces: Whittaker (1967, 1975)
% To cite two works by the same author with specific pages: \citeauthor{ref-journal-3a} (\citeyear{ref-journal-3a}, p. 328; \citeyear{ref-journal-3b}, p.475). This produces: Wong (1999, p. 328; 2000, p. 475)

%%%%%%%%%%%%%%%%%%%%%%%%%%%%%%%%%%%%%%%%%%
%% for journal Sci
%\reviewreports{\\
%Reviewer 1 comments and authors’ response\\
%Reviewer 2 comments and authors’ response\\
%Reviewer 3 comments and authors’ response
%}
%%%%%%%%%%%%%%%%%%%%%%%%%%%%%%%%%%%%%%%%%%
\PublishersNote{}
%\isPreprints{}{% This command is only used for ``preprints''.
\end{adjustwidth}
%} % If the paper is ``preprints'', please uncomment this parenthesis.
\end{document}

